% Irodov Part III §3.1, original 13--27. Source PDF pp.103--104,
% printed pp.94--95; answers PDF pp.305--306.
% Problems 17,18,25: textbook's printed coefficient conflicts with direct integration.
\electroproblem{13}{94}
一根长度为 $2a$ 的均匀带电细直杆总电荷为 $q$。
求距离杆中心为 $r$ 的点处电场模：
(a) 点在杆的中垂线上；
(b) 点在杆的延长线上（$r>a$）。
讨论 $r\gg a$。
\exsolution
设杆沿 $x$ 轴从 $-a$ 延伸至 $a$，
线电荷密度 $\lambda=q/(2a)$。
(a) 中垂线上点 $(0,r)$，横向分量抵消，
\[
E_\perp=k\lambda\int_{-a}^{a}
\frac{r\,\mathrm dx}{(r^2+x^2)^{3/2}}
=k\lambda\frac{2a}{r\sqrt{r^2+a^2}},
\]
因此
\[
\boxed{E_\perp=\frac{|q|}
{4\pi\varepsilon_0r\sqrt{r^2+a^2}}}.
\]
(b) 延长线上 $r>a$ 的点，
\[
E_\parallel=k\lambda\int_{-a}^{a}
\frac{\mathrm dx}{(r-x)^2}
=k\lambda\left(\frac1{r-a}-\frac1{r+a}\right),
\]
得到
\[
\boxed{E_\parallel=\frac{|q|}
{4\pi\varepsilon_0(r^2-a^2)}}.
\]
两种情况在 $r\gg a$ 时均有
$\boxed{E\simeq |q|/(4\pi\varepsilon_0r^2)}$。
\excheck 轴向表达式只在 $r>a$ 的杆外位置使用；
杆内理想线电荷的场需另外解释自场奇性。

\electroproblem{14}{94}
一根半无限长、均匀带电的直导线，
线密度为 $\lambda$，
求从其端点沿垂直方向距离为 $y$ 的点处
电场大小和方向。
\exsolution
取导线沿 $x\ge0$，观察点在 $(0,y)$。
电荷元 $\mathrm dq=\lambda\,\mathrm dx$，
位移向量为 $(-x,y)$。
因此
\[
E_x=-k\lambda\int_0^\infty
\frac{x\,\mathrm dx}{(x^2+y^2)^{3/2}}
=-\frac{k\lambda}{y},
\]
\[
E_y=k\lambda\int_0^\infty
\frac{y\,\mathrm dx}{(x^2+y^2)^{3/2}}
=\frac{k\lambda}{y}.
\]
合场大小
\[
\boxed{E=\frac{\sqrt2|\lambda|}
{4\pi\varepsilon_0y}}.
\]
\excheck 对正线电荷方向与直导线夹角为 $45^\circ$，
并从导线端点的上方指向远离导线的斜方向。

\electroproblem{15}{94}
一根很长、均匀线电荷密度 $\lambda$ 的细线
弯成图3.2的两种形状：
(a) 直角拐弯并由半径为 $R$ 的圆弧平滑连接；
(b) 半径为 $R$ 的对称 U 形弯曲。
求各图圆心 $O$ 的电场模；
直段长度远大于 $R$。
\exsolution
对每段细线用
\[
\mathrm d\bm E
=\frac{1}{4\pi\varepsilon_0}
\frac{\lambda\,\mathrm ds}{|\bm r_O-\bm r_s|^3}
(\bm r_O-\bm r_s)
\]
积分。长直段取半无限长极限。
(a) 直角弯的两直线段及圆弧产生的电场
按两个正交方向分解；
弧段的横、纵分量与直段部分抵消，
保留下来的两个分量大小均为 $k|\lambda|/R$。
故
\[
\boxed{E_a=\sqrt2\,\frac{k|\lambda|}{R}
=\frac{\sqrt2|\lambda|}{4\pi\varepsilon_0R}}.
\]
(b) U 形两侧直段的横向场分量因对称性相消，
其纵向和半圆弧的场积分大小相等、方向相反，
得到
\[
\boxed{E_b=0}.
\]
\excheck 结论与底本图3.2对应；
若把图(a)的两个圆弧中心误读为拐角点，
将得到错误的几何距离系数。

\electroproblem{16}{94}
半径为 $R$ 的球壳，表面电荷密度
$\sigma(\bm r)=\bm a\cdot\bm r$，
其中 $\bm a$ 为常矢量、
$\bm r$ 是从球心指向球面的矢量。
求球心处电场矢量。
\exsolution
令 $\bm r=R\bm n$，
$\mathrm dS=R^2\mathrm d\Omega$。
面元在球心产生的电场
\[
\mathrm d\bm E
=-k\frac{\sigma(\bm r)}{R^2}\bm n\,\mathrm dS
=-kR(\bm a\cdot\bm n)\bm n\,\mathrm d\Omega.
\]
利用球面角积分
$\int n_i n_j\,\mathrm d\Omega
=(4\pi/3)\delta_{ij}$，
得到
\[
\boxed{\bm E(0)=
-\frac{R}{3\varepsilon_0}\bm a}.
\]
\excheck 矢量方向与 $\bm a$ 相反，
大小和球壳半径 $R$ 成正比。

\electroproblem{17}{94}
半径为 $R$ 的球壳的面电荷分布为
$\sigma(\theta)=\sigma_0\cos\theta$，
$\theta$ 是对固定 $+z$ 轴的极角。
说明它可由两相反均匀带电实心球微小错位实现，
并求球壳内的电场矢量。
\exsolution
取两均匀体电荷密度为 $\pm\rho$ 的球
相对沿 $+z$ 轴错位 $\delta\ll R$。
两个球体表面相差的薄层在一阶近似中
形成面电荷密度
$\sigma(\theta)\simeq\rho\delta\cos\theta$，
故 $\rho\delta=\sigma_0$。
均匀带电实心球内部电场为
$\bm E=\rho(\bm r-\bm r_c)/(3\varepsilon_0)$。
正、负球的内部场叠加后，
$\bm r$ 项消掉，
\[
\bm E=
-\frac{\rho\delta}{3\varepsilon_0}\bm e_z
=\boxed{-\frac{\sigma_0}{3\varepsilon_0}\bm e_z}.
\]
\excheck \textbf{书末系数争议：}
本版简答印 $-\sigma_0\bm e_z/(2\varepsilon_0)$；
但直接球面积分
$-k\int\sigma_0\cos\theta\,\bm n\,\mathrm d\Omega$
和标准均匀极化球去极化场
均给系数 $1/3$，与题面所述球壳一致。
待核查书末是否混淆球壳与圆盘。

\electroproblem{18}{94}
半径为 $R$ 的实心球内，
体电荷密度
$\rho(\bm r)=\bm a\cdot\bm r$，
$\bm a$ 为常矢量。
求球心处电场。
\exsolution
取球坐标
$\bm r=r\bm n$，
体元
$\mathrm dV=r^2\mathrm dr\,\mathrm d\Omega$。
每一体元在球心的场为
\[
\mathrm d\bm E
=-k\frac{\rho(\bm r)}{r^2}\bm n\,\mathrm dV
=-k\,r(\bm a\cdot\bm n)\bm n\,
\mathrm dr\,\mathrm d\Omega.
\]
分离径向与角向积分：
\[
\bm E(0)
=-k\left(\int_0^Rr\,\mathrm dr\right)
\left(\int(\bm a\cdot\bm n)\bm n\,\mathrm d\Omega\right)
=-k\frac{R^2}{2}\frac{4\pi}{3}\bm a.
\]
故
\[
\boxed{\bm E(0)=
-\frac{R^2}{6\varepsilon_0}\bm a}.
\]
\excheck \textbf{书末系数争议：}
本版简答写 $-\bm aR^2/(2\varepsilon_0)$，
是上述体积分结果的3倍；
可由独立的对称张量平均
$\langle n_i n_j\rangle=\delta_{ij}/3$ 复核。

\electroproblem{19}{94}
一根半无限长均匀带电直线
沿半径为 $R$ 的圆盘轴线，
一端恰在圆盘中心，线密度为 $\lambda$。
求电场通过圆盘的通量，讨论法向选择。
\exsolution
在轴线上距离圆盘平面为 $z>0$ 的微元电荷
$\mathrm dq=\lambda\,\mathrm dz$，
圆盘张立体角大小为
$\Omega(z)=2\pi[1-z/\sqrt{z^2+R^2}]$。
对该点电荷，其电通量模为
$\mathrm d|\Phi|=(\mathrm dq/4\pi\varepsilon_0)\Omega(z)$。
于是
\[
|\Phi|=\frac{\lambda}{2\varepsilon_0}
\int_0^\infty
\left(1-\frac z{\sqrt{z^2+R^2}}\right)\mathrm dz.
\]
利用
$\int(1-z/\sqrt{z^2+R^2})\,\mathrm dz
=z-\sqrt{z^2+R^2}$，
在 $0,\infty$ 两端取值得 $R$，
因此
\[
\boxed{|\Phi|=\frac{|\lambda|R}{2\varepsilon_0}}.
\]
\excheck 圆盘是\emph{开曲面}，
通量符号取决于所选法向；
不能直接套闭曲面的高斯定理而忽视补面。

\electroproblem{20}{95}
图3.3中点电荷 $+q,-q$ 位于
半径为 $R$ 的圆盘两侧轴上，
分别距圆盘中心 $\ell$，
相距 $2\ell$。
求电场通过圆盘的通量。
\exsolution
每个点电荷在圆盘处张立体角
$\Omega=2\pi(1-\ell/\sqrt{\ell^2+R^2})$。
正电荷位于一侧，其场穿过圆盘；
负电荷位于另一侧，场线同向穿过，
两电荷贡献同号相加。
因此通量模为
\[
|\Phi|=\frac{2|q|}{4\pi\varepsilon_0}\Omega
=\boxed{\frac{|q|}{\varepsilon_0}
\left(1-\frac{\ell}{\sqrt{\ell^2+R^2}}\right)}.
\]
\excheck $R\to0$ 时通量趋零；
$R\to\infty$ 时通量模趋 $|q|/\varepsilon_0$；
正负号仍由圆盘法向定义决定。

\electroproblem{21}{95}
半径为 $R$ 的均匀体电荷密度
$\rho$ 的球体，
由距球心为 $r_0<R$ 的平面截出圆盘。
求电场穿过该截面的通量。
\exsolution
由均匀带电实心球的高斯定理，
球内位置矢量 $\bm r$ 的电场
$\bm E(\bm r)=\rho\bm r/(3\varepsilon_0)$。
取截面为 $z=r_0$，
截面内法向场分量处处相同：
$E_z=\rho r_0/(3\varepsilon_0)$。
圆盘半径
$b=\sqrt{R^2-r_0^2}$、
面积 $S=\pi(R^2-r_0^2)$。
因此
\[
\boxed{\Phi=
\frac{\pi\rho r_0(R^2-r_0^2)}{3\varepsilon_0}}.
\]
\excheck 若 $r_0=0$ 截面经过球心，
法向分量为零；若 $r_0=R$ 截面面积趋零，
通量同样趋零。

\electroproblem{22}{95}
两根很长的平行细线，
每根线电荷密度均为 $\lambda$，
相隔距离 $\ell$。
求位于两线中间的\emph{对称平面}内
电场强度的最大值。
\exsolution
令两线沿 $z$ 轴，
其横截面位置为
$x=\pm\ell/2,y=0$。
对称平面 $x=0$ 内，
两线产生的 $x$ 分量相消，
$y$ 分量相加：
\[
E_y(y)=
\frac{\lambda}{\pi\varepsilon_0}
\frac{y}{y^2+(\ell/2)^2}.
\]
令
$f(y)=|y|/[y^2+(\ell/2)^2]$，
对 $y>0$ 求导，
$f'(y)\propto(\ell/2)^2-y^2$。
极大出现在 $y=\ell/2$，故
\[
\boxed{E_{\max}=\frac{|\lambda|}{\pi\varepsilon_0\ell}}.
\]
\excheck 对称平面不是两线连线的平面；
沿二者中垂面 $x=0$，
恰在中点 $y=0$ 的场反而为零。

\electroproblem{23}{95}
无限长圆柱形薄壳上，
面电荷密度沿圆周方位角变化为
$\sigma(\varphi)=\sigma_0\cos\varphi$。
求圆柱轴线上的场强大小与方向。
\exsolution
取圆柱半径为 $R$，
观察点为轴线原点。
壳面元
$\mathrm dS=R\,\mathrm d\varphi\,\mathrm dz$，
电荷元
$\mathrm dq=\sigma_0\cos\varphi\,
R\,\mathrm d\varphi\,\mathrm dz$。
到轴线距离为
$\sqrt{R^2+z^2}$。
电场的横向 $\bm e_x$ 分量为
\[
E_x=-k\sigma_0 R^2
\int_0^{2\pi}\cos^2\varphi\,\mathrm d\varphi
\int_{-\infty}^{\infty}
\frac{\mathrm dz}{(R^2+z^2)^{3/2}}.
\]
两积分分别为 $\pi$ 和 $2/R^2$，
于是
\[
\boxed{\bm E=-\frac{\sigma_0}{2\varepsilon_0}\bm e_x}.
\]
\excheck 方向指向圆周上 $\varphi=\pi$ 的一侧，
场强与圆柱半径无关；
净线电荷密度沿圆周积分为零。

\electroproblem{24}{95}
某静电场只与 $x,y$ 有关：
\[
\bm E=a\,\frac{x\bm i+y\bm j}{x^2+y^2},
\]
其中 $a$ 常数。求通过球心位于原点、
半径为 $R$ 的球面的电通量。
\exsolution
在柱坐标中
$\bm E=(a/\varrho)\bm e_\varrho$，
是轴线附近具有奇性的径向场，
相当于等效无限线电荷场
$E_\varrho=\lambda/(2\pi\varepsilon_0\varrho)$。
匹配得 $\lambda=2\pi\varepsilon_0a$。
半径为 $R$ 的球面包围了轴线长为 $2R$
的一段有效线电荷，
$Q_{\rm enc}=\lambda(2R)$。
由高斯定律，
\[
\boxed{\Phi=Q_{\rm enc}/\varepsilon_0=4\pi aR}.
\]
\excheck 尽管 $\nabla\cdot\bm E=0$
在远离 $z$ 轴的点成立，
轴线处有分布意义的奇性；
不能由“处处散度零”误判球面通量为零。

\electroproblem{25}{95}
半径 $R$ 的绝缘球体电荷分布球对称：
$\rho(r)=\rho_0(1-r/R)$，$0\le r\le R$，
外部真空。求内外电场 $E(r)$、
最大值及其所处半径。
\exsolution
包围半径 $r\le R$ 内的总电荷为
\[
Q(r)=4\pi\int_0^r
\rho_0(1-s/R)s^2\,\mathrm ds
=4\pi\rho_0\left(\frac{r^3}3-\frac{r^4}{4R}\right).
\]
由高斯定律
$4\pi r^2E(r)=Q(r)/\varepsilon_0$，
\[
\boxed{E(r)=\frac{\rho_0r}{3\varepsilon_0}
\left(1-\frac{3r}{4R}\right)},\qquad r\le R.
\]
总电荷
$Q(R)=4\pi\rho_0R^3/12$，
所以外部
\[
\boxed{E(r)=\frac{\rho_0R^3}{12\varepsilon_0r^2}},
\qquad r\ge R.
\]
球内场求导
$E'(r)=\rho_0[1/3-r/(2R)]/\varepsilon_0$，
令其为零得
$r_{\max}=2R/3$，
\[
\boxed{E_{\max}=\frac{\rho_0R}{9\varepsilon_0}}.
\]
\excheck \textbf{本版书末第25(b)题系数疑误：}
简答写 $\rho_0R/(2\varepsilon_0)$，
代回其自身给出的 $E(r)$
在 $r=2R/3$ 处实际为
$\rho_0R/(9\varepsilon_0)$，后者由直接微分核实。

\electroproblem{26}{95}
半径为 $R$ 的球内部总电荷为 $q$，分布球对称；
球外介质从 $r=R$ 起另有空间电荷密度
$\rho(r)=\alpha/r$。
求让外部电场模与距离 $r$ 无关时，
球自身电荷 $q$ 的大小及该常场强。
\exsolution
由球对称的高斯定律，$r\ge R$ 时
\[
Q_{\rm enc}(r)
=q+4\pi\int_R^r(\alpha/s)s^2\,\mathrm ds
=q+2\pi\alpha(r^2-R^2).
\]
因此
\[
E(r)=\frac1{4\pi\varepsilon_0r^2}
[q-2\pi\alpha R^2+2\pi\alpha r^2].
\]
若 $E$ 与 $r$ 无关，
必须令 $1/r^2$ 的系数为零：
\[
\boxed{q=2\pi\alpha R^2},\qquad
\boxed{E=\frac{\alpha}{2\varepsilon_0}}.
\]
\excheck 电荷密度虽随 $1/r$ 减小，
球外包围的附加电荷仍按 $r^2$ 增长，
恰能维持常量场。

\electroproblem{27}{95}
空间中存在球对称电荷密度
$\rho(r)=\rho_0e^{-\alpha r^3}$，
$\rho_0,\alpha>0$。
求 $E(r)$，并研究
$\alpha r^3\ll1$ 与 $\alpha r^3\gg1$。
\exsolution
\textbf{题面图像复核：} 指数为
$-\alpha r^3$，PDF 文字抽取常丢失立方上标。
由高斯定律
\[
E(r)=\frac{1}{\varepsilon_0r^2}
\int_0^r\rho_0s^2e^{-\alpha s^3}\,\mathrm ds.
\]
令 $u=\alpha s^3$，则
$s^2\mathrm ds=\mathrm du/(3\alpha)$，
故
\[
\boxed{E(r)=\frac{\rho_0}
{3\varepsilon_0\alpha r^2}
(1-e^{-\alpha r^3})}.
\]
当 $\alpha r^3\ll1$，
$1-e^{-u}\simeq u$，所以
\[
\boxed{E(r)\simeq\frac{\rho_0r}{3\varepsilon_0}}.
\]
当 $\alpha r^3\gg1$，
指数项可略，得到
\[
\boxed{E(r)\simeq\frac{\rho_0}
{3\varepsilon_0\alpha r^2}}.
\]
\excheck 近球心呈均匀带电球内的线性规律；
远处呈有限总电荷产生的 $1/r^2$ 规律，
两个极限均正确。
